\section{n3. Positive-dimensional intersections of distinct irreducible components obstruct geometric envelopes in every characteristic.}
\textit{Supports: G331, n1, n2.}

A general cross-component conductor obstruction follows from the ENTIRE qualified G331, without assuming that the normal components have CE. Let R be any excellent reduced positive-dimensional local ring, a a primary gauge, S its finite normalization, rho its FULL boundary quotient and c its conductor. Suppose one rho-fibre has lifts alpha,beta on normal local centres belonging to DIFFERENT minimal components of R. Then some bounded strictly negative continuous SIGNED PRIMARY-MODEL obstacle on R has neither a greatest geometric C nor a greatest geometric D minorant.

Choose a regular c0 in the conductor. The normalization has idempotents e1,e2 for the two distinct minimal-domain factors. Since the conductor is an S-ideal, j=c0 e1 and ell=c0 e2 belong to R. They vanish identically on every other component and are nonzero on their own domain factors. The nontrivial fibre lies on the infinite-conductor stratum by G331, so alpha,beta kill c0. These are actual elements downstairs, not hypothetical separated branch functions.

Choose a centred divisorial valuation nu on the normal local centre containing beta, normalized by aS. G331's normal-domain divisorial bound supplies a finite K\textgreater{}=1 such that nu(I)\textless{}=K beta(I) for every unit-or-primary ideal on this centre. Precisely, the nonunit bound nu(h)\textless{}=C ord\_m(h), together with beta(m)\textgreater{}=1/r for m\textasciicircum{}r subset aS, gives nu(I)\textless{}=Cr beta(I); units are separate. Consequently every negative finite REAL primary combination satisfies u(nu)\textgreater{}=K u(beta). Uniform comparison and actual decreasing infima give this for ENTIRE geometric C and D. Pullback of a target g therefore yields g(rho(beta))\textless{}=g(rho(nu))/K. This uses exactly G331's normal bound, its actual cone comparison and full-support pullback. It introduces no analytic irreducibility claim for nonnormal R.

The number d=nu(ell)=nu(c0) is positive finite. Choose integers L\textgreater{}d and B\textgreater{}0 so that 1+B(1-d/L)\textgreater{}K. Put b=B+1 and the ACTUAL proper primary ideal I=(ell)+a\textasciicircum{}L. Define the homogeneous signed primary-model obstacle
F=(B/L)v(I)-(B+1)v(a).
On V\_a it lies in[-b,-1]; on component1 it is exactly-1, and its value at the common fibre is-1. At rho(nu) it is -1-B(1-d/L)\textless{}-K. It is continuous by primary-ideal evaluation and is a signed REAL Cartier evaluation on a directed common primary blowup. Every geometric minorant g\textless{}=F therefore obeys g(rho(beta))\textless{}=F(rho(nu))/K\textless{}-1.

For every n\textgreater{}=1 put T\_n=a\textasciicircum{}(bn)+(j)a\textasciicircum{}n and q\_n=-v(T\_n)/n. These are actual proper primary ideals. Their negative effective REAL Cartier divisors on ACTUAL Bl\_(T\_n) are relatively ample, hence have positive degree on EVERY contracted integral curve. On component1 the profile is max(-b,-1-v(j)/n), and on every other component it is-b. Thus q\_n\textless{}=F on EVERY support. At each component1 point with finite j-value their supremum is-1. This locus is dense in the full component1 boundary: G331's all-support divisorial-density theorem gives a dense subset of generic-support centred divisorial valuations, and j is nonzero there. In particular alpha lies in its closure, even though alpha kills the conductor.

A greatest C or D minorant would dominate every q\_n and thus equal-1 on this dense locus (F itself is-1 there). Its normalization pullback is USC. If its value at alpha were\textless{}-1, an open strict sublevel neighbourhood of alpha would meet the dense locus, a contradiction. But alpha and beta have the same target point, where EVERY minorant is strictly below-1 by the preceding inequality. Therefore no greatest C/D minorant exists. This is a proof using dense nets in general; no metrizability or sequence extraction is assumed.

An algebraic sufficient condition for the fibre hypothesis is that two distinct minimal primes p1,p2 satisfy dim R/(p1+p2)\textgreater{}=1. Take a centred semivaluation of the reduction of this excellent positive-dimensional quotient, whose boundary is nonempty by G331. Its support contains both primes. For each i its factor through R/p\_i lifts by finite normalization lifting to a normal centre of the i-th normalization domain. These give two lifts in different components of the same full rho-fibre. Their conductor value must be infinite by uniqueness on the finite stratum. Hence EVERY such positive-dimensional intersection of distinct irreducible components obstructs the all-continuous-obstacle envelope property, in ANY dimension, characteristic or residue field. The produced obstacle is already a signed primary model.

This does not address nonsingleton fibres lying entirely within a single minimal-domain normalization, nor all excellent normal dimension\textgreater{}=3 envelopes. No arbitrary-model cofinality, resolution, degree averaging, source CE, newly defined ideal cone, or omitted conductor support is used. Every statement is gauge covariant: F and q\_n are homogeneous, positive gauge transport divides by the same bounded denominator. Artifacts n3.md and n3.tex; rerun python deliver.py; python check\_delivery.py. The universal geometric obstruction is mathematical reasoning relative to G331, not a finite computation or an original-task closure.

\subsection*{Cross-component signed model obstruction}
For actual conductor elements $j=c_0e_1$, $\ell=c_0e_2$ choose
$L>\nu(\ell)$ and $B$ with $1+B(1-\nu(\ell)/L)>K$. Then
\[
 F=\frac BLv((\ell)+\mathfrak a^L)-(B+1)v(\mathfrak a),\qquad
 q_n=-\frac{v(\mathfrak a^{(B+1)n}+(j)\mathfrak a^n)}n.
\]
\[
 g(\rho\beta)\le\frac{g(\rho\nu)}K<-1,\qquad
 \sup_n q_n=-1\text{ on a dense first-component locus}.
\]
This contradicts upper semicontinuity of any greatest minorant at
$\rho\alpha=\rho\beta$; the constants use the granted NORMAL bound.

